\chapter[Formal Two-Law Schema]{\titlesplit{The Two Laws}{Formal Axioms and Schema}}\label{app:two-laws}

This appendix fixes the technical contract used by the later appendices. The
Constitutive Law determines the admissible extension cone. The revised
Selective Law has three ordered faces: \PEN{} selects a structural API inside
that cone; free sealing gives the least lawful realization of the selected API;
and physical public actualization requires stable internal records satisfying
Public Confluence. The public-event extractor itself remains bridge content
unless further observer dynamics derive it.

All definitions below are relative to a fixed ambient CCHM-style cubical
extension calculus \(\mathfrak{C}\), equipped with cumulative universes,
\(\Pi\)- and \(\Sigma\)-types, finite sums, path types, univalence, Kan
composition, normalized public signatures, and sealed opaque extensions. The
necessity-level two-law schema is presentation-independent; the particular
audit clauses in this appendix are the book's present canonical specification.
Extraction-Guarded Debt--Credit Provenance, introduced below, is an explicit
strengthening of that specification. It is not claimed to follow from the
earlier Constitutive and Selective axioms.

\section{State, Signatures, and Candidates}

\begin{definition}[History and public boundary]
A \emph{history} \(H_n\) is a finite sequence of sealed extensions
\[
H_n = (x_1,\dots,x_n)
\]
interpreted in \(\mathfrak{C}\), modulo definitional equality and univalent
equivalence of normalized public signatures. Its exported interface is the
normalized telescope
\[
\partial H_n := \mathrm{NF}\bigl(\partial x_1,\dots,\partial x_n\bigr).
\]
Write \(\mathcal{L}(H_n)\) for the finite audited basis of derivation schemas
generated by \(\partial H_n\) under the fixed closure operator
\(\mathrm{Cl}_{\mathfrak{C}}\). Equivalent schemas are counted once.
\end{definition}

\begin{definition}[Candidate continuation]
A candidate continuation over \(H\) is a finite normalized extension telescope
\[
x = (\partial x,\; w_x,\; O_\bullet(x\mid H),\; b_x),
\]
where \(\partial x\) is the proposed public boundary, \(w_x\) is its constructive
realization datum, \(O_\bullet(x\mid H)\) is the normalized obligation complex
induced by sealing \(x\) against \(\partial H\), and
\[
b_x : \partial H \longrightarrow \partial(H,x)
\]
is the boundary inclusion preserving the already sealed public interface.
\end{definition}

\begin{definition}[Presentation equivalence]
Two candidates \(x,y\) over \(H\) are presentation-equivalent, written
\[
x \simeq_H y,
\]
when their sealed public signatures are connected by a univalent equivalence
over \(\partial H\) preserving normalized kernel clauses, computation clauses,
and the audited derivation-schema basis. All quantities below are required to be
invariant under \(\simeq_H\).
\end{definition}

\begin{definition}[Grammar and expressible cone]
A grammar \(G\) is a finite normalized constructor signature. The raw
expressible cone \(\mathrm{Cand}_G(H)\) is the set of normalized candidates over
\(H\) buildable from \(G\) and the public constructors already exported by
\(\partial H\). The admissible cone will be a CL-filtered, horizon-bounded
subcone of \(\mathrm{Cand}_G(H)\).
\end{definition}

\begin{definition}[Free sealing completion]
Let \(\mathcal{S}_x\) denote the normalized generators, recursors, computation
rules, and coherence equations of the selected candidate \(x\). Its free
sealing completion over \(H\) is
\[
F_x(H)
:=
\mathrm{Term}_{\mathfrak C}(H;\mathcal{S}_x)/{\equiv_x},
\]
where \(\mathrm{Term}_{\mathfrak C}(H;\mathcal{S}_x)\) is generated from the
old library and \(\mathcal{S}_x\), and \(\equiv_x\) is the least congruence,
stable under substitution, composition, and the exported public operations,
containing the old judgmental equalities and the equations forced by
\(\mathcal{S}_x\). No additional generator or equation is included.
\end{definition}

\begin{definition}[Integration]
If \(x\) is selected at \(H\), the integration map seals its free completion
into history:
\[
I(H,x) := \mathrm{Seal}_{\mathfrak{C}}\bigl(F_x(H)\bigr),
\qquad
H^+ = I(H,x).
\]
Integration normalizes the public boundary, hides nonpublic construction data
behind the sealed extension, and exports only the charged public API clauses
and their definitional consequences.
\end{definition}

\begin{theorem}[Relative free-sealing initiality]
\label{thm:relative-free-sealing}
Assume the generated reduction system of \(F_x(H)\) is lawful in
\(\mathfrak C\). Let \(\mathbf{Comp}_x(H)\) be the category of lawful exact
realizations of \(\mathcal{S}_x\) extending \(H\), with interpretations over
\(H\) that preserve the selected operations and forced equations. Then
\(F_x(H)\) is initial in \(\mathbf{Comp}_x(H)\):
\[
\mathrm{Hom}_{H,x}\bigl(F_x(H),C\bigr)\cong\mathbf 1
\qquad
\text{for every }C\in\mathbf{Comp}_x(H).
\]
When the recursion, induction, \(\beta\)-, \(\eta\)-, and coherence data are
retained intensionally, the corresponding mapping type is contractible.
\end{theorem}

\begin{proof}
Interpret inherited generators by the given map from \(H\), interpret each new
generator by its specified realization in \(C\), and extend recursively over
generated terms. The interpretation descends through \(\equiv_x\) because
\(C\) satisfies every forced equation. Any second interpretation agrees on the
old library and on every selected generator, hence agrees on every generated
term by structural induction and on every equivalence class by quotient
induction. The intensional version uses the retained \(\eta\)- and coherence
rules to upgrade uniqueness to contractibility.
\end{proof}

\paragraph{What initiality does not say.}
The theorem is relative to the API already selected. It does not say that PEN
selects the initial completion of an independently specified weaker
requirement. Minimal positive overshoot and categorical leastness are different
orders: a richer API can lie closer to the bar than a free presentation of a
weaker requirement. The universal property governs realization after
selection; it is not a second proof of the selection ranking.

\section{The Constitutive Law}

The Constitutive Law is the predicate \(\mathrm{CL}_H(x)\) determining which
candidates can coherently count as possible continuations of \(H\).

\begin{definition}[Constitutive admissibility]
\(\mathrm{CL}_H(x)\) holds iff the following four clauses hold.
\begin{enumerate}[leftmargin=*, label=\textbf{CL\arabic*.}]
\item \textbf{Witness-bearing realizability.}
\[
\partial H,\partial x \vdash_{\mathfrak{C}} \mathrm{ok}
\quad\text{and}\quad
\partial H \vdash_{\mathfrak{C}} w_x : \mathrm{Realize}(x\mid H).
\]
Existence claims are admitted only with constructive realization data.

\item \textbf{Equivalence invariance.} If \(x\simeq_H y\), then
\[
\mathrm{CL}_H(x)\Longleftrightarrow \mathrm{CL}_H(y),
\qquad
I(H,x)\simeq I(H,y).
\]
Ontology is counted at the level of normalized public structure, not syntactic
redescription.

\item \textbf{Bounded local satisfiability.} The obligation complex
\(O_\bullet(x\mid H)\) has finite normalized support, and every primitive
obligation is discharged inside the active constitutive depth bound
\[
a(o) := |\mathrm{Supp}(o)| \le d_{\mathrm{obl}},
\qquad o\in O_\bullet(x\mid H).
\]
In the canonical cubical realization used in this book,
\(d_{\mathrm{obl}}=2\); the proof of exactness is deferred to
\hyperref[app:d2-proof]{Appendix \ref*{app:d2-proof}}.

\item \textbf{Canonicity under sealing.} Every closed exported construction in
\(\partial I(H,x)\) normalizes to a canonical public value, path, computation
rule, or opaque constant; no sealed candidate may export a stuck neutral as a
primitive public commitment.
\end{enumerate}
\end{definition}

\begin{theorem}[Constitutive Law]
A candidate can enter realized history only if it is witness-bearing,
equivalence-invariant, locally satisfiable at bounded constitutive depth, and
canonical after sealing:
\[
x\in\mathcal{A}_{G,h}(H)\quad\Rightarrow\quad \mathrm{CL}_H(x).
\]
\end{theorem}

\section{Audited Novelty and Cost}

\begin{definition}[Kernel clauses and integration cost]
The kernel of \(x\) over \(H\) is the minimal normalized public clause set
\[
\mathrm{KernelClauses}(x\mid H)
\]
such that
\[
\mathrm{Cl}_{\mathfrak{C}}\bigl(\partial H\cup
\mathrm{KernelClauses}(x\mid H)\bigr)
\simeq
\mathrm{Cl}_{\mathfrak{C}}\bigl(\partial I(H,x)\bigr)
\]
over \(\partial H\). The integration cost is the discrete audit count
\[
\kappa_H(x)
=
\bigl|\mathrm{KernelClauses}(x\mid H)\bigr|
\in \mathbb{N}.
\]
Definitional completion, aliases, projections forced by already charged
clauses, and presentation changes are not charged. If \(\kappa_H(x)=0\), then
\(x\) is a theorem-readout or redescription of \(H\), not a candidate for a new
realized step.
\end{definition}

\begin{definition}[Typed operational schema audit]
\label{def:typed-operational-schema-audit}
For a history \(H\), let \(\mathsf{RawSch}_2(H)\) be the well-typed raw
derivation-schema judgements in the fixed depth-two closure fragment of
\(\mathfrak C\). A judgement records its parameter context, conclusion, and
derivation; an untyped expression is not a schema. Write \(\approx_H\) for the
least equivalence relation generated by judgmental equality and by transport
along univalent equivalences of normalized public signatures, subject to
preservation of the conclusion and of all substitution and naturality squares.

An operational audit supplies a terminating typed normalizer
\[
  \mathrm{nf}_H:\mathsf{RawSch}_2(H)\longrightarrow
  \mathsf{NfSch}_2(H)
\]
with certificates of type preservation, soundness, and completeness:
\[
  \mathrm{nf}_H(D)=\mathrm{nf}_H(E)
  \quad\Longleftrightarrow\quad
  D\approx_H E.
\]
Full normalization of all of \(\mathfrak C\) is not required; a complete
normalizer for the finite auditable depth-two fragment suffices. If these
certificates are unavailable, the schema audit fails closed rather than
replacing semantic equality by equality of printed syntax.
\end{definition}

\begin{definition}[Schema family and instance]
\label{def:schema-family-instance}
A semantic schema family over \(H\) is a typed normal schema varying naturally
in its declared parameters: reindexing a parameter context by any well-typed
substitution agrees, after normalization, with substituting into the family.
Two families are equal when they are pointwise \(\approx_H\)-equal. Let
\[
  \mathsf{Sch}_2(H)
\]
be the finite set of these family classes in the audited fragment.

For a family \(\sigma\) and a well-typed substitution \(\theta\), the term
\(\sigma[\theta]\) is an \emph{instance} of \(\sigma\), not a second family.
Specializing one universe-polymorphic family at \(N\) sealed types therefore
does not by itself contribute \(N\) units of novelty. A specialization becomes
a separate counting unit only when the public boundary exports it as an
independently typed family and proves it inequivalent to the other exported
families, or when it fills a distinct independently exported demand orbit as
defined below. In either case the audit promotes it to an orbit-indexed
zero-parameter family before forming \(\mathsf{Sch}_2(H)\); raw instances never
enter that set directly.
\end{definition}

\begin{definition}[Weakening and marginality]
\label{def:schema-weakening-marginality}
The boundary inclusion \(b_x:\partial H\to\partial I(H,x)\) induces typed
normalized weakening
\[
  \mathrm{wk}_x:\mathsf{Sch}_2(H)\longrightarrow
  \mathsf{Sch}_2(I(H,x)).
\]
The marginal semantic families of \(x\) are
\[
  \operatorname{Marg}_2(H;x)
  :=
  \mathsf{Sch}_2(I(H,x))\setminus\operatorname{im}(\mathrm{wk}_x).
\]
Thus an old family remains old after substitution, renaming, or univalent
transport. A derived term using a new parameter is marginal only when its
normalized family class does not lie in the weakening image. This definition
separates the semantic quotient from any later rule deciding which marginal
families receive positive credit.
\end{definition}

\begin{definition}[Individual semantic demand orbits]
\label{def:semantic-demand-orbits}
A \emph{demand instance} of a depth-two window is a typed normalized required
output together with its support, origin stage, and discharge judgement. Two
instances occupy the same semantic demand orbit exactly when context
reindexing, definitional equality, natural-family substitution, or univalent
transport identifies both the required output and its discharge judgement.
Write
\[
  \mathsf{LiveOrb}(H)
\]
for the finite set of orbits generated before the candidate under audit and not
derivable in \(H\). Each orbit \(o\) carries a finite set
\(\mathsf{Req}_H(o)\) of independently required output positions.

An orbit-ledger certificate must prove \emph{extraction completeness}---it
enumerates every normalized demand of the active window---and
\emph{derivability completeness}---it decides, with a typed derivation or a
live witness, the disposition of every extracted orbit. It must also carry a
\emph{window-locality and expiration} proof: every demand that can be live at
stage \(n+1\) is generated by the active depth-two window
\((S_n,S_{n-1})\), while every demand of an older window was discharged and
remains discharged under cumulative weakening. Without this third clause,
emptying the current window does not imply that the whole live-orbit set is
empty. A focus-family name may annotate an orbit, but equality or emptiness of
focus labels is not evidence for equality or emptiness of semantic demand
orbits.
\end{definition}

\begin{proposition}[Orbit-level sufficiency for the J2/J3 debt check]
\label{prop:orbit-ledger-j2-j3}
For the active Step-16 window \((S_{15},S_{14})\), an
extraction-complete orbit inventory together with a
derivability-complete proof that every orbit is discharged in \(B_{15}\), and
a depth-two window-locality/expiration proof that these are all orbits that can
be live at stage~16, is sufficient for the uses of J2 totality and J3 closure
in the debt-exhaustion argument. It follows that the semantic directive debt
at Step~16 is empty.
This certificate does not derive J2's free initiality or J3's law-like kernel
from A1--A4; those remain additional hypotheses unless proved separately.
\end{proposition}

\begin{proof}
Extraction completeness rules out an omitted demand instance of the
\((S_{15},S_{14})\) window. Derivability completeness places every resulting
orbit in the normalized derivability closure of \(B_{15}\). Window locality
rules out a live orbit sourced anywhere else, and expiration plus monotone
weakening prevents an older discharged orbit from returning. The set difference
defining semantic directive debt therefore has no element. This establishes
exactly the total-discharge and closure consequences used in the debt-exhaustion
proof, while saying nothing by itself about the universal property of the
completion or why its eight clauses are law-like.
\end{proof}

\paragraph{Status of the orbit refinement.}
Proposition~\ref{prop:orbit-ledger-j2-j3} is a sufficiency criterion, not a
report that its certificate has already been constructed. The existing
focus-family replay establishes emptiness only at package-label granularity.
Until a typed extractor enumerates every demand orbit of
\((S_{15},S_{14})\), decides each orbit modulo weakening and univalent
equality, and proves the depth-two window-locality/expiration clause, semantic
\(\mathsf{LiveOrb}(H_{15})=\varnothing\) remains an explicit proof obligation.

\begin{definition}[Extraction-Guarded Debt--Credit Provenance axiom]
\label{ax:extraction-guarded-provenance}
Let
\[
  \mathsf K_H(x):=\mathrm{KernelClauses}(x\mid H)
\]
and let the strengthened law's local role code be the following four-element
inductive set:
\[
  \mathcal R_2=\{
  \mathsf{KernelHead},
  \mathsf{AdjointMate},
  \mathsf{SupportAction},
  \mathsf{Coherence}\}.
\]
The roles mean respectively: the irreducible public declaration; its one
adjoint or eliminative face; one natural P5, P6, or synthesis support action;
and one independent depth-two coherence or interchange family. Define the
finite provenance-tag type
\[
  \mathsf{ProvTag}_H(x)
  :=
  \bigl(\mathsf K_H(x)\times\mathcal R_2\bigr)
  \;\sqcup\;
  \coprod_{o\in\mathsf{LiveOrb}(H)}\mathsf{Req}_H(o).
\]

For \(f\in\operatorname{Marg}_2(H;x)\) and
\(t\in\mathsf{ProvTag}_H(x)\), the proposition
\(\mathsf{Anchors}_{H,x}(f,t)\) has the following two constructors.
\begin{enumerate}[leftmargin=*, label=\textbf{A\arabic*.}]
\item If \(t=(a,r)\) lies in the local summand, an anchor is a typed
normal-form proof that \(a\in\mathsf K_H(x)\) is the first irreducible fresh
source of \(f\), together with a derivation that \(f\) actually has local role
\(r\).
\item If \(t=(o,q)\) lies in the demand summand, an anchor is a typed
normal-form discharge proof showing that the pre-existing live orbit \(o\) and
required output \(q\in\mathsf{Req}_H(o)\) are realized by \(f\).
\end{enumerate}
In particular, a string label, a cardinal slot, or membership in the same
focus package is not an anchor.

The strengthened Selective Law requires every counted audit to carry an
equivalence-invariant injection
\[
  \pi_{H,x}:\operatorname{Marg}_2(H;x)\hookrightarrow
  \mathsf{ProvTag}_H(x)
  \tag{EGP}
\]
together with the dependent witness
\[
  \alpha_{H,x}:
  \prod_{f\in\operatorname{Marg}_2(H;x)}
  \mathsf{Anchors}_{H,x}\bigl(f,\pi_{H,x}(f)\bigr).
  \tag{Anchor}
\]
Write \(\mathrm{EGP}_H(x)\) when both \(\pi_{H,x}\) and
\(\alpha_{H,x}\) exist. A demand first created by \(x\) cannot occur in the
demand summand: its independent public content must instead be charged in
\(\mathsf K_H(x)\).

The role grammar, orbit equivalence, and injection are \emph{bar-blind} and
stage-generic. They may not depend on \(n\), \(\mathrm{Bar}_n\), the desired
winner, the accepted trace, or a survivor list. They must commute with typed
substitution, weakening, and univalent presentation equivalence. Without an
EGP certificate, \(\nu_H^{\mathrm{tot}}(x)\) is undefined for Selective-Law
purposes and \(x\) cannot be ranked.

This is an explicit additional axiom of the strengthened Selective Law. It is
motivated by earned novelty and anti-packing, but it is not a theorem of the
Constitutive Law, the prior five PEN rules, free sealing, or univalence alone.
In particular, existence of \((\pi_{H,x},\alpha_{H,x})\) is the required
anchored four-role exhaustion and thinness certificate. Merely listing the four
constructors or injecting into four anonymous slots does not prove that a
candidate in the intended semantics admits such a certificate.
\end{definition}

\begin{definition}[Generative novelty]
For an EGP-certified operational audit, the total novelty of \(x\) over \(H\)
is the number of marginal natural-family classes
\[
\nu^{\mathrm{tot}}_H(x)
=
\bigl|\operatorname{Marg}_2(H;x)\bigr|.
\]
Equivalently, this is
\(\bigl|\mathcal{L}(I(H,x))\setminus\operatorname{im}(\mathrm{wk}_x)\bigr|\)
when \(\mathcal L\) is realized by the typed operational schema basis above.
Writing a bare set difference \(\mathcal L(I(H,x))\setminus\mathcal L(H)\) is
only shorthand; weakening is the map that identifies the old basis inside the
new one.
The canonical audit decomposes this as
\[
\nu^{\mathrm{tot}}_H(x)
=
\nu^G_H(x)+\nu^C_H(x)+\nu^K_H(x),
\]
where \(\nu^G\) counts formation/introduction schemas, \(\nu^C\) counts
elimination/projection/cross-interface schemas, and \(\nu^K\) counts
computation, path, Kan, and coherence-reduction schemas.
\end{definition}

\begin{definition}[Efficiency ratio]
For \(\kappa_H(x)>0\) and a complete EGP-certified operational audit, the
efficient novelty ratio is
\[
\rho_H(x)=\frac{\nu^{\mathrm{tot}}_H(x)}{\kappa_H(x)}.
\]
The ratio is undefined for \(\kappa_H(x)=0\), and such candidates are excluded
from the continuation cone.
\end{definition}

\begin{proposition}[Provenance accounting bound]
\label{prop:provenance-accounting-bound}
Every EGP-certified candidate satisfies
\[
  \bigl|\operatorname{Marg}_2(H;x)\bigr|
  \le
  |\mathcal R_2|\,\kappa_H(x)
  +
  \sum_{o\in\mathsf{LiveOrb}(H)}|\mathsf{Req}_H(o)|.
  \tag{P}
\]
Consequently, on a semantically debt-free field, any independently proved
bar-blind bound \(|\mathcal R_2|\le c\) gives
\[
  \nu_H^{\mathrm{tot}}(x)\le c\,\kappa_H(x),
  \qquad
  \rho_H(x)\le c.
\]
For the four-constructor code fixed by
Axiom~\ref{ax:extraction-guarded-provenance}, this specializes to
\[
  \nu_H^{\mathrm{tot}}(x)\le4\kappa_H(x),
  \qquad
  \rho_H(x)\le4.
\]
\end{proposition}

\begin{proof}
Take finite cardinalities in the injection (EGP), use
\(|\mathsf K_H(x)|=\kappa_H(x)\), and distribute cardinality over the disjoint
union. If \(\mathsf{LiveOrb}(H)=\varnothing\), the second summand is empty.
The four role constructors are pairwise distinct, so
\(|\mathcal R_2|=4\).
\end{proof}

\begin{definition}[Certified Step-15-generated guarded flow]
\label{def:guarded-step15-internal-flow}
Let \(H_{15}\) be the sealed result of the Step-15 temporal/modal completion.
A \emph{Step-15-generated guarded flow} \(g\) is a post-completion continuation
over that sealed API whose declared operational content is transport or
coiteration through the guarded structure already exported by Step~15. It is
not the Step-15 candidate itself. Such a flow is \emph{certifiably internal}
when
its typed schema audit supplies an erasure map
\[
  \mathrm{erase}_g:\mathsf{Sch}_2(I(H_{15},g))\longrightarrow
  \mathsf{Sch}_2(H_{15})
\]
natural under substitution and univalent transport, and the following inverse
laws hold on family classes:
\[
  \mathrm{erase}_g\circ\mathrm{wk}_g=\mathrm{id},
  \qquad
  \mathrm{wk}_g\circ\mathrm{erase}_g=\mathrm{id}.
  \tag{WE}
\]
Operationally, erasure removes the guarded Step-15 flow parameters and
normalizes the transported derivation in the sealed old library. Merely
asserting that a carrier is equivalent to an old carrier is not an erasure
certificate; the maps must preserve the exported operations, computation
rules, and naturality data.
\end{definition}

\begin{theorem}[Conservative internality for Step-15-generated flows]
\label{thm:guarded-step15-internality}
Every certifiably internal guarded continuation \(g\) has
\[
  \operatorname{Marg}_2(H_{15};g)=\varnothing,
  \qquad
  \nu^{\mathrm{tot}}_{H_{15}}(g)=0.
\]
\end{theorem}

\begin{proof}
For any \(\tau\in\mathsf{Sch}_2(I(H_{15},g))\), the second inverse law gives
\[
  \tau=\mathrm{wk}_g(\mathrm{erase}_g(\tau)).
\]
Thus weakening is surjective. Its image is the whole target, so the complement
defining \(\operatorname{Marg}_2\) is empty and its finite cardinality is zero.
The unique map and dependent witness from the empty set supply
\((\pi,\alpha)\), so the audited novelty is defined and equals zero.
\end{proof}

\paragraph{Status of the internality theorem.}
The argument above proves internality from the weakening/erasure inverse laws.
It does not yet construct \(\mathrm{erase}_g\), or prove (WE), for every proposed
post-completion flow generated by the Step-15 API. Those are typed
normalization obligations, not consequences of merely labelling a continuation
internal. The theorem makes no claim that the Step-15 completion itself had
zero novelty.

\begin{proposition}[Blind finite local-role bound]
\label{prop:blind-four-role-bound}
The local-role coefficient of the strengthened Selective Law is independent of
the stage, bar, history, and survivor set, and is exactly
\[
  |\mathcal R_2|=4.
\]
Every semantically debt-free EGP-certified candidate therefore satisfies
\[
  \nu_H^{\mathrm{tot}}(x)\le4\kappa_H(x),
  \qquad
  \rho_H(x)\le4.
\]
\end{proposition}

\begin{proof}
The four inductive constructors in
Axiom~\ref{ax:extraction-guarded-provenance} are distinct and the definition
consults none of the listed historical quantities. The EGP injection enforces
at most one marginal family per kernel-clause/role tag, while
\(\alpha\) proves that every such tag is an actual semantic anchor. Apply
Proposition~\ref{prop:provenance-accounting-bound} with an empty live-orbit set.
\end{proof}

\paragraph{Logical status of the four-role result.}
This is a theorem of the \emph{strengthened} Selective Law because an EGP
certificate is, by definition, a faithful injection into the four-role code.
It is not a theorem of the previous Two-Law axioms. The remaining semantic work
is to construct such a certificate for every candidate to be ranked, and to
replay the historical stages. If that construction needs a larger code, the
strengthening must be revised openly; any independently justified code of at
most nine roles would retain the weaker \(9\kappa\) ceiling.

\begin{corollary}[Conditional Step-16 exclusion]
\label{cor:conditional-step16-egp-exclusion}
Assume all of the following:
\begin{enumerate}[leftmargin=*]
\item an extraction-, derivability-, and window-locality-complete orbit
certificate proves \(\mathsf{LiveOrb}(H_{15})=\varnothing\);
\item the provenance re-audit of Steps \(1\)--\(15\) preserves the selected
history \(H_{15}\): at every stage it certifies the revised candidate cone,
selective order, and winner, retains the fifteen audited scores, and hence
retains
\[
  \mathrm{Bar}_{16}=\frac{354333}{39040};
\]
\item the Step-16 cone is exhausted by certifiably internal guarded flows and
opaque candidates carrying EGP certificates.
\end{enumerate}
Then no Step-16 candidate clears the bar. Every opaque candidate has
\[
  \rho_{H_{15}}(x)\le4
  <
  \frac{354333}{39040},
\]
while every certifiably internal flow has \(\nu=0\).
\end{corollary}

\begin{proof}
The empty orbit ledger removes the demand summand from (EGP), so
Proposition~\ref{prop:blind-four-role-bound} bounds every opaque candidate by
four. The strict comparison is \(4\cdot39040=156160<354333\).
Theorem~\ref{thm:guarded-step15-internality} handles the internal branch. The
exhaustive dichotomy then covers the Step-16 cone.
\end{proof}

\paragraph{Status of the exclusion corollary.}
This corollary is not yet an unconditional global-halt theorem. It still awaits
the typed demand-orbit certificate, the provenance-based re-audit of all
fifteen historical candidate cones, orders, winner certificates, and scores,
weakening/erasure inverse witnesses for the intended internal flows, and EGP
certificates establishing the stated Step-16 exhaustion.

\paragraph{Audit conventions.}
All counts are taken after normalization and quotienting by
\(\simeq_H\). A public API clause is charged exactly once, at the first stage at
which it is irreducible over the current history. A derived eliminator is free
when it is definitionally forced by an already charged adjoint or universal
property; it is charged when it contributes a new public observation,
projection, or cross-interface rule not reconstructible from the prior public
boundary.

\paragraph{Reconciliation of legacy amplification rules.}
The traditional schema-calculus labels P5 and P6, the \(d^2\) and \(r^2\)
formulas, and combinatorial schema synthesis remain useful descriptions of raw
capacity. Under Axiom~\ref{ax:extraction-guarded-provenance}, they are not
automatic novelty multipliers.

\begin{enumerate}[leftmargin=*, label=\textbf{R\arabic*.}]
\item \textbf{P5 record inheritance.} A lifted inherited family in the
weakening image is old. A genuinely marginal lift earns independent credit
only through a separately paid kernel-role tag or a distinct pre-existing
demand output. Importing a large historical surface does not by itself copy
that surface into the numerator.

\item \textbf{P6 uniform action.} One polymorphic natural family is counted
once. Its specializations at library entries are not multiplied unless they
are independently exported, inequivalent family classes and carry distinct
provenance tags---in particular, distinct demand-orbit outputs when the kernel
does not separately pay for them.

\item \textbf{Dimension-squared and reference-squared capacity.} The ordered
\((i,j)\) filling or transport sites described by \(d^2\) or \(r^2\) may measure
raw expressibility. Each site counted as independent novelty must itself be a
distinct normalized natural family and must either be separately charged or
answer a pre-existing live demand. Sites forced by one family, or obtained by
specialization and composition, are instances or definitional completion.

\item \textbf{Combinatorial synthesis.} A lower bound on distinct raw
\(\Pi\)-/\(\Sigma\)-expressions is not yet a lower bound on marginal semantic
families. The synthesis multiplier survives only after typed normalization
proves the families inequivalent and the EGP injection assigns each of them a
distinct paid or demanded source.
\end{enumerate}

Independent semantic credit must therefore be paid independently or discharge
an independently live demand. This rule reconciles amplification with the
Selective Law's earned-novelty philosophy; it does not deny the corresponding
constructions exist in the Constitutive cone.

For the Genesis audit, a minimal API bundle \(\mathcal{S}\) importing
\(k\) prior interfaces with unique maximal imported layer \(L_{\max}\) uses the
telescopic capacity rule
\[
\nu^C_H(\mathcal{S})
=
\nu^{\mathrm{tot}}_H(L_{\max})
+ \kappa_H(\mathcal{S}) + (k-1).
\]
This is an audit rule for counted public capacity, not an additional
constitutive axiom. Under the strengthened Selective Law it is valid only when
the inherited term, local faces, and bridges possess the EGP provenance just
specified; the equation is not by itself their certificate.

\paragraph{Historical re-audit obligation.}
Axiom~\ref{ax:extraction-guarded-provenance} applies at every stage, including
the stages whose scores motivated it. All fifteen Genesis values must
therefore be recomputed from typed marginal families and individual demand
orbits, and selection must be rerun on the revised certified candidate cone at
each stage. An earlier amplification may be retained when the replay exhibits the
pre-existing live demand outputs or separately charged roles that supported it.
Numerical replay of the old formula without that correspondence is a regression
check, not a semantic proof; preserving the scores without preserving the
candidate cones, ordering, and winner certificates does not preserve the
Genesis trace. Until this orbit-level re-audit is complete, the historical
ledger under the strengthened law is provisional.

\section{The Selective Law}

The Selective Law is the transition rule that chooses the next realized
continuation from the CL-admissible cone.

\begin{definition}[Active horizon and admissible cone]
The active cost horizon at stage \(n\) is \(h_n\in\mathbb{N}_{>0}\). Given
history \(H_n\) and grammar \(G_n\),
\[
\mathcal{A}_{G_n,h_n}(H_n)
=
\{\,x\in\mathrm{Cand}_{G_n}(H_n):
\mathrm{CL}_{H_n}(x),\;
0<\kappa_{H_n}(x)\le h_n\,\}.
\]
After a successful realization, \(h_{n+1}=2\). After an idle search tick with no
bar-clearing prime candidate, \(h_n\) is replaced by \(h_n+1\).
\end{definition}

\begin{definition}[Rising bar]
Let \(\nu_i=\nu^{\mathrm{tot}}_{H_{i-1}}(x_i)\) and
\(\kappa_i=\kappa_{H_{i-1}}(x_i)\) be the audited values of the already realized
steps. Let \(\Delta_i\) be the presentation-invariant integration debt of layer
\(i\) in the \(d=2\) lane. For \(n>1\), define
\[
\Omega_{n-1}
=
\frac{\sum_{i<n}\nu_i}{\sum_{i<n}\kappa_i},
\qquad
\Phi_n
=
\frac{\Delta_n}{\Delta_{n-1}},
\qquad
\mathrm{Bar}_n
=
\Phi_n\,\Omega_{n-1}.
\]
The initial seed convention for \(n=1\) is fixed by the starting state. A
candidate clears the bar at stage \(n\) iff
\[
\rho_{H_n}(x)\ge \mathrm{Bar}_n.
\]
\end{definition}

\begin{definition}[Constructive primeness]
Let \(\Gamma_x\) be the finite dependency graph of the normalized public clauses
of \(x\), with an edge \(a\to b\) when \(b\) occurs irreducibly in the type,
computation rule, or coherence boundary of \(a\). A nonempty proper union
\(T\subsetneq\Gamma_x\) of terminal strongly connected components determines a
subcandidate \(x|T\) when its public boundary seals over \(H\).

The candidate \(x\) is \emph{composite} at stage \(n\) iff some such \(T\)
satisfies
\[
x|T\in\mathcal{A}_{G_n,h_n}(H_n),
\qquad
\rho_{H_n}(x|T)\ge \mathrm{Bar}_n.
\]
It is \emph{constructively prime}, written \(P_{H_n}(x)\), iff it is not
composite.
\end{definition}

\begin{definition}[Selective ordering]
Among candidates satisfying
\[
x\in\mathcal{A}_{G_n,h_n}(H_n),
\qquad
\mathrm{EGP}_{H_n}(x),
\qquad
\rho_{H_n}(x)\ge\mathrm{Bar}_n,
\qquad
P_{H_n}(x),
\]
the Selective Law uses the lexicographic order
\[
x\prec_n y
\quad\Longleftrightarrow\quad
\bigl(
\rho_{H_n}(x)-\mathrm{Bar}_n,\;
\kappa_{H_n}(x),\;
\mathrm{key}_{H_n}(x)
\bigr)
<
_{\mathrm{lex}}
\bigl(
\rho_{H_n}(y)-\mathrm{Bar}_n,\;
\kappa_{H_n}(y),\;
\mathrm{key}_{H_n}(y)
\bigr),
\]
where \(\mathrm{key}_{H_n}\) is a deterministic normalized-presentation
tie-breaker. Equivalently, one may use the penalty objective
\[
\mathcal{J}_n(x)
=
\mathcal{P}_{\mathrm{adm}}(x)
+\mathcal{P}_{\mathrm{prov}}(x)
+\mathcal{P}_{\mathrm{bar}}(x)
+\mathcal{P}_{\mathrm{prime}}(x)
+\gamma\bigl(\rho_{H_n}(x)-\mathrm{Bar}_n\bigr)
+\epsilon\kappa_{H_n}(x),
\]
where the four penalties are \(0\) respectively on admissible, EGP-certified,
bar-clearing, and prime candidates, and \(+\infty\) otherwise, with
\(\gamma>0\) and
\(\epsilon>0\) chosen so that the finite part realizes the same lexicographic
priority.
\end{definition}

\begin{theorem}[Selective Law]
Assume the normalized horizon-bounded cone is finite, and let
\[
\mathcal{C}_n
=
\{\,x\in\mathcal{A}_{G_n,h_n}(H_n):
\mathrm{EGP}_{H_n}(x),\;
\rho_{H_n}(x)\ge\mathrm{Bar}_n,\;
P_{H_n}(x)\,\}.
\]
If \(\mathcal{C}_n\neq\varnothing\), actuality selects
\[
x_{n+1}
=
\min_{\prec_n}\mathcal{C}_n
\]
and integrates
\[
H_{n+1}=I(H_n,x_{n+1}).
\]
The selected extension is unique up to \(\simeq_{H_n}\) after applying the
normalized key. This is uniqueness in the PEN ranking on the current viable
frontier. It is not categorical initiality among completions of a weaker
requirement. The categorical statement is instead
Theorem~\ref{thm:relative-free-sealing}, applied after \(x_{n+1}\) has been
selected.
\end{theorem}

\section{PEN as the Operational Encoding}

The \PEN{} is the following executable transition schema for the Selective Law.
At stage \(n\), the state is
\[
S_n=(H_n,G_n,h_n,\Delta_n,\Omega_n,\mathsf{LiveOrb}(H_n)).
\]

\begin{enumerate}[leftmargin=*, label=\textbf{PEN\arabic*.}]
\item \textbf{Enumerate.} Enumerate normalized anonymous candidates in
\(\mathrm{Cand}_{G_n}(H_n)\) whose kernel cost is within the active horizon.

\item \textbf{Constitute.} Type-check candidates in \(\mathfrak{C}\), normalize
their public signatures, quotient by \(\simeq_{H_n}\), and retain exactly
\(\mathcal{A}_{G_n,h_n}(H_n)\).

\item \textbf{Audit.} Produce the typed normal forms, natural-family quotient,
weakening map, individual demand-orbit inventory, and EGP injection. Then
compute
\[
\operatorname{Marg}_2(H_n;x),\qquad
\pi_{H_n,x},\qquad
\kappa_{H_n}(x),\qquad
\nu^G_{H_n}(x),\qquad
\nu^C_{H_n}(x),\qquad
\nu^K_{H_n}(x),\qquad
\rho_{H_n}(x).
\]

\item \textbf{Prime-filter.} Remove candidates composite by terminal-SCC
subbundle decomposition.

\item \textbf{Select by least sufficient overshoot.} Select the unique
normalized minimizer under \(\prec_n\) among candidates with
\(\rho_{H_n}(x)\ge\mathrm{Bar}_n\).

\item \textbf{Integrate freely.} Form the relative free completion of the
selected API and seal its normalized public interface:
\[
H_{n+1}=I(H_n,x_{n+1})
=\mathrm{Seal}_{\mathfrak C}\bigl(F_{x_{n+1}}(H_n)\bigr),\qquad
G_{n+1}=\mathrm{NF}(G_n\cup\partial x_{n+1}),
\qquad
h_{n+1}=2.
\]
Update \(\Delta_{n+1}\), \(\Omega_{n+1}\), the semantic demand-orbit ledger,
and the next bar from the audited ledger.
\end{enumerate}

If PEN1--PEN4 produce no bar-clearing prime candidate, no history step is
realized; the search performs the idle update \(h_n\mapsto h_n+1\) and repeats.

\section{Physical Sealing and Public Confluence}

The preceding clauses determine formal library growth. A physical public event
requires an internal record interface in addition.

\begin{definition}[Internal public seal]
For an internal context \(p\), a public seal for a formally integrated event
\(x\) consists of a record carrier \(R_p\), a distinguishable record state
\(r_x:R_p\), and admissible comparison maps to prior and overlapping public
records. The seal is \emph{stable} when it is repeatable, copyable through those
comparison maps, and preserved by later admissible evolution until replaced by
an explicitly declared sufficient statistic proved lossless for every retained
public claim.
\end{definition}

\begin{definition}[Internal Record Actualization axiom]
\label{ax:internal-record-actualization}
A formal integration becomes a public physical event for \(p\) only when it has
an internal public seal that is Born-supported by the physical dynamics,
decodable to the declared tolerance, compatible with prior lock-in, and inside
the active record-capacity bound. Formal selection without such a seal is not
yet a public event for \(p\).
\end{definition}

\begin{definition}[Public-event extractor]
A public-event extractor for context \(p\) is a partial map
\[
\mathcal Q_p:
(H_{\mathrm{pub}},\rho_{pW},x)
\rightharpoonup
\mathsf{Rec}_p
\]
whose defined values satisfy Axiom~\ref{ax:internal-record-actualization}. If
the conditions fail, \(\mathcal Q_p\) returns no public seal rather than a
contradictory one.
\end{definition}

\begin{definition}[Public Confluence axiom]
\label{ax:public-confluence}
Let \(r_p\) and \(r_q\) be separately stable public seals concerning a common
overlap. If they are operationally compatible, there exists a joint
nondemolition refinement \(r_{pq}\), unique up to public operational
equivalence, with comparison maps
\[
r_{pq}\longmapsto r_p,
\qquad
r_{pq}\longmapsto r_q
\]
that preserve both seals. If no such refinement exists, no realized public
history contains both as facts about that overlap. Finite-capacity failure is
therefore ``common seal'' or ``no common seal,'' never two mutually
contradictory public seals.
\end{definition}

\begin{proposition}[Extractor and kinematic underdetermination]\label{prop:extractor-underdetermination}
The Constitutive Law, PEN, relative free sealing, Internal Record Actualization,
and Public Confluence do not uniquely determine \(\mathcal Q_p\), nor do they
select a unique nonpublic operational state-space geometry.
\end{proposition}

\begin{proof}
Distinct record POVMs, tolerances, and coarse-grainings can satisfy the same
stability and confluence clauses while leaving every structural PEN quantity
unchanged. Likewise, classical simplex theories, real and quaternionic quantum
fragments, and generalized probabilistic models can carry the same confluent
classical public-record interface. Public Confluence constrains the common
record layer; further reconstruction principles are required to determine the
private state cone, pure-process category, and composite rule.
\end{proof}

\begin{proposition}[Model-relative extractor derivability]\label{prop:model-relative-extractor}
Universal extractor underdetermination is compatible with exact
observer-relative derivation.  There exists a finite calibrated observer model
in which the dynamics uniquely determines the minimal sufficient public
statistic for its retained event claim, the optimal mandatory-decision error
floor, the append-only record instrument, every least common refinement in its
public algebra, and the pointer and noise resource entropies.
\end{proposition}

\begin{proof}
\hyperref[app:finite-public-extractor]{Appendix~\ref*{app:finite-public-extractor}}
constructs the model and proves each assertion.  For a balanced binary event
copied through \(N\) independent thermal flip cells, likelihood-equivalence
classes are exactly the Hamming-weight classes.  Their projectors define the
instrument; strict likelihood ordering fixes the optimal no-abstention
decision; intersections of their coarse-grained partitions give exact
nondemolition gluing; and collision probabilities give the resource
entropies.  The construction depends on its declared observer dynamics and
therefore does not contradict Proposition~\ref{prop:extractor-underdetermination}.
\end{proof}

\section{Anti-Packing and Scope}

The cost model excludes axiom-packing by construction. An opaque constant with
no public eliminator, computation rule, or cross-interface action contributes
bounded \(\nu\) and cannot acquire arbitrary novelty merely by being named. A
structured bundle that exports several independent bar-clearing submoves is
composite by the terminal-SCC criterion. A bundle whose clauses are mutually
irreducible is not composite, but then every irreducible public commitment is
charged in \(\kappa\); packing cannot lower cost by hiding public obligations in
private syntax.

The necessity-level architecture is therefore:
\[
\begin{aligned}
\text{Constitutive Law}
&:\quad
\mathrm{Cand}_G(H)\longrightarrow \mathcal{A}_{G,h}(H),\\
\text{PEN selection}
&:\quad
\mathcal{A}_{G,h}(H)\longrightarrow x^\star,\\
\text{free integration}
&:\quad
x^\star\longrightarrow H^+=I(H,x^\star),\\
\text{public physical sealing}
&:\quad
(H^+,p)\mathrel{\mathop{\dashrightarrow}^{\mathcal Q_p}}H^+_{\mathrm{pub}},
\qquad H^+_{\mathrm{pub}}\text{ satisfies Public Confluence}.
\end{aligned}
\]
The canonical PEN implementation sharpens this with the audited quantities
\(\nu\), \(\kappa\), \(\rho\), \(\mathrm{Bar}\), and constructive primeness.
Relative free sealing supplies initiality only after the API is selected;
Internal Record Actualization and Public Confluence govern the additional
physical-public face.

\hyperref[app:genesis-table]{Appendix~\ref*{app:genesis-table}} records the strict Genesis trace under this contract. \hyperref[app:genesis-script]{Appendix~\ref*{app:genesis-script}}
states the Genesis script in generative form. \hyperref[app:d2-proof]{Appendix~\ref*{app:d2-proof}} proves the exact
\(d_{\mathrm{obl}}=2\) depth result used by the bar dynamics. \hyperref[app:minimal-extension]{Appendix~\ref*{app:minimal-extension}} applies
the same contract to the minimal grammar-extension theorem at the univalent
horizon.
